\section{Arithmetic localization of the missing carrier}
\label{sec:tpc60-missing-partition}

The previous sections are exact for any pair of finite-dimensional maps.
We now return to the fixed-\(h_0\) Walsh carrier and record what the literal
arithmetic dictionary does and does not add.

\subsection{Boolean signatures and ordered first failure}

For \(w\in\cW_H(y)\), recover \(m,j,\gamma,n,p,r,s\) by
\eqref{eq:tpc60-target-n}--\eqref{eq:tpc60-prs-recovery}.  Membership in
the represented terminal image is a finite conjunction of declared tests.
For definiteness, use the dependency order
\begin{align}
 T_1(w)&:=\one_{\{(m,r)\in\cI_X\}},
 \label{eq:tpc60-test-retained-pair}\\
 T_2(w)&:=\one_{\{j,\gamma,p\text{ lie in the declared terminal scopes}\}},
 \label{eq:tpc60-test-terminal-scopes}\\
 T_3(w)&:=\one_{\{\Lambda_X(n)\ne0\}},
 \label{eq:tpc60-test-terminal-weight}\\
 T_4(w)&:=\one_{\{d_m,r,s\text{ squarefree and }(d_m,r)=1\}},
 \label{eq:tpc60-test-squarefree}\\
 T_5(w)&:=\one_{\{\text{all remaining terminal-only masks pass}\}}.
 \label{eq:tpc60-test-remaining-masks}
\end{align}
The lists inside \(T_2\) and \(T_5\) are the inherited literal declarations,
not new favorable supports.  By construction,
\begin{equation}
 w\in\cW_R
 \quad\Longleftrightarrow\quad
 T_1(w)=\cdots=T_5(w)=1.
 \label{eq:tpc60-represented-all-tests}
\end{equation}

The Boolean signature \((T_1,\ldots,T_5)\) is order independent.  A
``first failure'' is not: it becomes well defined only after the order
above is declared.  Define
\begin{align}
 \cF_1&:=\{w\in\cW_H:T_1(w)=0\},
 \label{eq:tpc60-first-failure-1}\\
 \cF_\nu&:=\{w\in\cW_H:
 T_1(w)=\cdots=T_{\nu-1}(w)=1, T_\nu(w)=0\},
 \quad2\le\nu\le5.
 \label{eq:tpc60-first-failure-nu}
\end{align}

\begin{proposition}[Exact ordered missing partition]
\label{prop:tpc60-exact-missing-partition}
The first-failure sets are disjoint and
\begin{equation}
 \boxed{
 \cW_M=\bigsqcup_{\nu=1}^5\cF_\nu.}
 \label{eq:tpc60-missing-first-failure-partition}
\end{equation}
If
\begin{equation}
 \cD_\nu(\zeta):=\sum_{w\in\cF_\nu}|b_w(\zeta)|^2,
 \label{eq:tpc60-first-failure-energy}
\end{equation}
then
\begin{equation}
 \boxed{
 \cD_M=\sum_{\nu=1}^5\cD_\nu.}
 \label{eq:tpc60-DM-failure-partition}
\end{equation}
\end{proposition}

\begin{proof}
Every missing atom fails at least one test by
\eqref{eq:tpc60-represented-all-tests}; its least failing index is unique.
Conversely every atom in a first-failure set fails terminal admission and
is missing.  Raw atomic additivity gives the energy identity.
\end{proof}

The adjective ``ordered'' is essential.  Reordering the tests changes the
first-failure labels, although it does not change \(\cW_M\) or the full
Boolean signature.

\subsection{Empty pseudo-failure layers}

Several apparent failure mechanisms are already excluded on the common
active Walsh carrier.

\begin{proposition}[Inherited automatic tests]
\label{prop:tpc60-inherited-automatic-tests}
For an actual nonzero high Walsh atom on the common TPC--45/TPC--53
physical carrier:
\begin{enumerate}[label=\textup{\(\alph*\)},leftmargin=2.5em]
 \item \(p=P^+(n)>y\), \(pr=mj+h_0\), and \(r=n/p\) hold by the unique
       high-prime decoding;
 \item the nonzero Walsh amplitude already enforces
       \(\Lambda_X(n)\ne0\) and its shared literal profile masks;
 \item the inherited squarefree and coprimality support implies that
       \(r\) and \(s=d_mr\) are squarefree and \((d_m,r)=1\).
\end{enumerate}
Consequently the first-failure layers created solely by \(T_3\) and \(T_4\)
are empty.  The part of \(T_5\) consisting only of masks shared with the
Walsh carrier is empty as well.
\end{proposition}

\begin{proof}
Part \(a\) is the largest-prime census and the assumption
\(y>\sqrt{N_+}\).  For \(b\), coordinates with zero shaped weight or zero
shared profile were removed as identically inactive when the common Walsh
carrier was formed.  For \(c\), the TPC--45 support makes \(d_mn\)
squarefree and \((d_m,n)=1\).  Since \(n=pr\), every divisor \(r\) is
squarefree, is coprime to \(d_m\), and therefore \(s=d_mr\) is squarefree.
\end{proof}

\begin{corollary}[Identical-support specialization]
\label{cor:tpc60-identical-support-missing}
Suppose the Walsh and terminal constructions use identical row, orbit,
opposite-row, terminal-prime, literal-mask, and profile supports, and there
is no additional terminal-only pruning.  Then
\begin{equation}
 \boxed{
 \cW_M
 =\{w\in\cW_H(y):(m(w),r(w))\notin\cI_X\}.}
 \label{eq:tpc60-missing-only-cofactor}
\end{equation}
Thus the only missing-high mechanism in this specialization is failure of
the complementary cofactor to enter the retained-pair carrier.
\end{corollary}

This is a support-alignment theorem, not an occupancy theorem.  It does not
show that the set in \eqref{eq:tpc60-missing-only-cofactor} has small mass,
or even that some active row contains a represented atom.

\subsection{Zero-cost localization of the diagnostic branch}

Let \(q\le5\) be the number of nonempty first-failure layers, and write
\begin{equation}
 N_{\rm dyad}(y)
 :=1+\left\lfloor\log_2\frac{N_+}{y}\right\rfloor
 =O(\log X)
 \label{eq:tpc60-dyadic-block-count}
\end{equation}
for the number of represented cofactor cells.  If
\(\cD_M>0\), one of them satisfies
\begin{equation}
 \max_\nu\cD_\nu\ge\frac{\cD_M}{q}.
 \label{eq:tpc60-failure-pigeonhole}
\end{equation}
Combining this constant partition with the \(O(\log X)\) dyadic partition
of the represented terminal face gives the exact dichotomy
\begin{equation}
 \cD_H^W:=\cD_R+\cD_M>0
 \quad\Longrightarrow\quad
 \begin{cases}
  \exists k:\displaystyle
  \cD_k\ge\dfrac{\cD_H^W}{2N_{\rm dyad}(y)},
  &\cD_R\ge\cD_H^W/2,\\[8pt]
  \exists\nu:\displaystyle
  \cD_\nu\ge\dfrac{\cD_H^W}{2q},
  &\cD_M\ge\cD_H^W/2,
 \end{cases}
 \label{eq:tpc60-represented-missing-dichotomy}
\end{equation}
Hence localization has zero fixed-power cost in
either branch:
\begin{equation}
 \boxed{\lambda_{\rm part}=0.}
 \label{eq:tpc60-zero-partition-cost}
\end{equation}
The conclusion is diagnostic.  It either finds a large represented block
or identifies a dominant reason for missing mass.  It does not force the
represented alternative and does not imply
\(\lambda_{\rm reasm}=0\).

\subsection{Exact rowwise represented-mass minimax}

Write
\begin{equation}
 W_E(m):=\sum_{\substack{w\in\cW_E\\m(w)=m}}|\beta_w|^2,
 \qquad E\in\{R,M,H\},
 \qquad W_H:=W_R+W_M.
 \label{eq:tpc60-row-weights}
\end{equation}
Then literal factorization gives
\begin{equation}
 \cD_E(\zeta)=\sum_mW_E(m)|\zeta_m|^2.
 \label{eq:tpc60-row-diagonal-form}
\end{equation}

\begin{proposition}[Sharp represented-mass minimax]
\label{prop:tpc60-rowwise-represented-minimax}
The optimal constant uniform over any bounded coefficient cube containing
a nonzero multiple of each coordinate direction is
\begin{equation}
 \boxed{
 \inf_{\cD_H^W(\zeta)>0}
 \frac{\cD_R(\zeta)}{\cD_H^W(\zeta)}
 =\min_{m:W_H(m)>0}\frac{W_R(m)}{W_H(m)}.}
 \label{eq:tpc60-rowwise-represented-minimax}
\end{equation}
In particular, if one active row has
\(W_R(m)=0<W_H(m)\), the uniform represented-mass transfer is exactly zero.
\end{proposition}

\begin{proof}
The quotient is a weighted average of the row ratios
\(W_R(m)/W_H(m)\), with nonnegative weights
\(W_H(m)|\zeta_m|^2\).  It is at least their minimum.  Concentrating
\(\zeta\) on a minimizing row attains that minimum after a harmless scalar
rescaling into the bounded cube.
\end{proof}

Thus a geometric density statement about the union of rows cannot replace
a rowwise lower bound when the coefficient class has only upper bounds on
\(|\zeta_m|\).

\subsection{Residual ladders and ambient degree}

Fix a common source-prime cell of diameter \(O(L)\), a residual key
\((i,s)\), and a dyadic cofactor block \(R\le r<2R\).  Write
\begin{equation}
 m=d\ell,
 \qquad r=\frac sd.
 \label{eq:tpc60-d-ell-r}
\end{equation}
The fixed-shift equality becomes
\begin{equation}
 rp-dj\ell=h_0.
 \label{eq:tpc60-determinant-equation}
\end{equation}

\begin{proposition}[Ambient residual ladder bound]
\label{prop:tpc60-ambient-residual-degree}
On one common \(D,L\) cell, the number of full, represented, missing, or
first-failure high atoms with fixed residual key \((i,s)\) and
\(R\le r<2R\) satisfies
\begin{equation}
 \boxed{
 d_{\max}^{\rm res}(R)
 \ll_{h_0}\tau(s)\left(1+\frac LR\right)
 =X^{o(1)}\left(1+\frac LR\right).}
 \label{eq:tpc60-residual-degree-bound}
\end{equation}
Moreover the prime-resolved degrees at \((i,p)\) and \((i,s,p)\) are at
most one.  The direct native degree has only the inherited crude bound
\begin{equation}
 d_{\rm miss}\ll X^{o(1)}Q.
 \label{eq:tpc60-crude-native-degree}
\end{equation}
\end{proposition}

\begin{proof}
Fix \(d\mid s\), hence \(r=s/d\), and put
\(g=(r,dj)=(r,j)\), using \((d,r)=1\).  If
\eqref{eq:tpc60-determinant-equation} has one integral solution
\((p_0,\ell_0)\), all integral solutions lie on the determinant ladder
\begin{equation}
 p=p_0+\frac{dj}{g}t,
 \qquad
 \ell=\ell_0+\frac rg t,
 \qquad t\in\mathbb Z.
 \label{eq:tpc60-determinant-ladder}
\end{equation}
Solvability forces \(g\mid h_0\), and therefore \(g\le|h_0|\).  Since
\(\ell\) lies in an interval of diameter \(O(L)\), the ladder contributes
\(O_{h_0}(1+L/R)\) candidates.  Summing over \(d\mid s\) gives
\eqref{eq:tpc60-residual-degree-bound}; primality, retained-pair, residue,
  mask, and profile tests only delete candidates.  The fixed-
  \((i,s,p)\) assertion is the TPC--55/TPC--58 prime-resolved injection,
  while the stronger fixed-\((i,p)\) assertion is the TPC--59
  critical-cutoff row singleton proved again in
  \cref{prop:tpc60-same-prime-separation}.

  For the native bound, fixing \(i\) fixes the side, opposite row, and
  \(j\).  Fixed-coordinate fused-label injectivity leaves at most one
  active high atom per source row, and the inherited source-system support
  bound is
  \(\#\mathcal A_\sigma\ll X^{o(1)}Q\), \(\sigma\in\{1,2\}\).
  This proves \eqref{eq:tpc60-crude-native-degree}.  A wider union of
  \(D,L\) cells must pay explicitly for the additional cells.
\end{proof}

The residual estimate is not the native estimate used directly in
\eqref{eq:tpc60-mass-degree-upper}: erasing \(s\) can merge many residual
fibers.  Prime-resolved singleton structure therefore does not imply
native singleton structure.
